\section{Introduction}
\label{sec:introduction}

Post-harvest agricultural losses present a major challenge to global food security and supply chain sustainability \cite{fao2019state}. The Food and Agriculture Organization (FAO) estimates that approximately 14\% of global food production is lost between harvest and retail distribution \cite{fao2019state}, with commercial fruit losses consistently ranging between 7\% and 16\% (e.g., 10.39\% in apple, 7.76\% in banana, and 15.88\% in citrus) \cite{jha2015assessment}. Manual visual sorting is labor-intensive, subjective, and prone to inspector fatigue \cite{rizzo2023fruit,zhang2024fruit}. While convolutional neural networks (CNNs) such as MobileNetV2 \cite{sandler2018mobilenetv2} have achieved over 95\% accuracy on curated fruit-quality benchmarks \cite{chen2024classification,gulzar2023fruit}, standard softmax classifiers suffer from severe overconfidence: normalized probability distributions peak sharply near 1.0 even on novel species or corrupted imagery \cite{guo2017calibration,hendrycks2017baseline}. In automated packing facilities, an overconfident false positive---such as misrouting an unseen species or a chemically adulterated lot into a fresh output stream---creates severe commercial and food-quality liabilities.

Post-hoc out-of-distribution (OOD) detection algorithms aim to identify and reject anomalous inputs at inference time without requiring model retraining \cite{scheirer2013toward,bendale2016towards,yang2024generalized}. Methods such as Maximum Softmax Probability (MSP) \cite{hendrycks2017baseline}, input-perturbed temperature scaling (ODIN) \cite{liang2018enhancing}, Mahalanobis feature distance \cite{lee2018simple}, and Energy-based detection \cite{liu2020energy} have demonstrated success on synthetic benchmarks \cite{vaze2022open}. However, existing agricultural evaluations conflate distribution shifts: when anomalies are drawn from external repositories while in-distribution data originates from curated laboratory archives, detector performance is confounded by differences in camera hardware, ambient illumination, and background environments rather than true botanical or pathological divergence.

To resolve these limitations, this study conducts a source-stratified empirical benchmark evaluating a frozen MobileNetV2 six-class fruit-quality classifier across $N = 13,465$ photographic records from three independent repositories: Kaggle Fruit Quality, FruitVision \cite{bijoy2025fruitvision}, and Sultana et al.\ \cite{sultana2022extensive}. By enforcing source-matched comparisons, we isolate external domain shift from semantic species and dataset-reported chemical treatment shifts, addressing three core research questions:
\begin{itemize}
\setlength{\itemsep}{0pt}\setlength{\parsep}{0pt}\setlength{\parskip}{0pt}\setlength{\topsep}{1pt}\setlength{\partopsep}{0pt}
\item \textbf{RQ1 (External Domain Robustness):} How severely does closed-set accuracy degrade when transferring from laboratory benchmarks to external photographic repositories, and what are the specific class-wise failure modes?
\item \textbf{RQ2 (Comparative OOD Discrimination):} How effectively do post-hoc scoring functions separate in-distribution produce from unseen botanical species versus dataset-reported chemical preservative shifts under source-matched conditions?
\item \textbf{RQ3 (Operational Gating Feasibility and Trade-offs):} What operational rejection and false-acceptance trade-offs emerge under fixed, validation-calibrated operating thresholds?
\end{itemize}

\noindent\textbf{Contributions:} (1) A source-stratified evaluation framework across five partitions and four cohorts spanning $13,465$ images. (2) Diagnostic quantification of a 28.56 percentage-point closed-set accuracy drop on external data, driven by an asymmetric collapse on Fresh Apples (32.12\% accuracy). (3) Evaluation of five post-hoc OOD algorithms and an operational Energy--MSP Dual-Gate cascade, demonstrating moderate discrimination on semantic species shifts (AUROC up to 81.57\%) but failure on dataset-reported chemical treatment shifts (AUROC $\le 58.77\%$, Dual-Gate UFAR-Fresh $= 42.83\%$). (4) Perceptual-hash sensitivity analysis establishing that detector rankings remain strictly invariant across deduplication tiers.
